\subsection{Simultaneous open Frey images}
\label{subsec:frey-open-images}
\paragraph{A rank-one weight obstruction.}

For $\ell=2,3$, the geometric image of the Legendre Tate module is
prosolvable: $\operatorname{SL}_2(\F_\ell)$ is solvable and the
successive congruence kernels are abelian. Thus both representations
are defined on the quotient under consideration.

\begin{lemma}[Common pure Frobenius polynomials]\label{lem:frey-common-frobenius}
For a noncuspidal section in Theorem~\ref{thm:prosolvable-frey-open-intro},
there is a Dirichlet-density-one set of rational primes at which both
Frey representations are unramified and have the same polynomial
$T^2-a_qT+q$, $a_q\in\Z$, with complex roots of modulus $q^{1/2}$.
\end{lemma}
\begin{proof}
Proposition~\ref{prop:frey-three-density} makes the local representative
interior with integral Legendre parameter away from the discriminant
on a density-one set, after removing the fixed finite bad-model set.
Both representations are then the Tate modules of the same elliptic
curve with good reduction at $q$. Its Frobenius polynomial and Weil
bound give the assertion.
\end{proof}

\begin{lemma}\label{lem:frey-rank-one-weights}
Let \(\psi:G_L\to E^\times\) be a continuous character of a number field,
where \(E/\Q_\ell\) is finite. If it is Hodge--Tate at all places over
\(\ell\), then it has finite ramification and is the \(\ell\)-adic
realization of an algebraic character of a Serre group
\(S_{\mathfrak m}\), defined over a number field.
It has realizations at every coefficient prime.
For some integer \(b\), at every place \(v\) outside its modulus and
different from the coefficient prime its algebraic Frobenius value has
absolute value \((Nv)^{b/2}\) under every complex embedding.
\end{lemma}
\begin{proof}
Tate's local theorem gives local algebraicity:
\cite[III \S1.2, III-7]{Serre1989}.
Finite ramification for abelian continuous representations, and the
existence of a modulus, are proved in
\cite[III \S2.2, III-10--11]{Serre1989}.
These statements are only being applied to a character.

There is a rationality qualification in the stated global theorem
\cite[III \S2.3, Theorem~2, III-13]{Serre1989}.
Its proof, III-14--15, first constructs a morphism from
\(S_{\mathfrak m}\) over the coefficient field, and uses rationality only
in the final descent to \(\Q\).
In the rank-one case that first morphism is a character after extending
coefficients to \(\overline\Q_\ell\).
Character groups of groups of multiplicative type do not acquire new
elements upon extending algebraically closed fields: the group-like
elements of a diagonalizable group's coordinate group algebra are its
character monomials. Hence this character descends through
\(\overline\Q\hookrightarrow\overline\Q_\ell\) to a finite number field.

Composition with the canonical representation
\(\varepsilon_{\ell'}:G_L^{\mathrm{ab}}\to
S_{\mathfrak m}(\Q_{\ell'})\) gives each companion. The common algebraic
Frobenius values are supplied by
\cite[II \S2.5, proposition, II-18--19]{Serre1989}.
The absolute-value formula is exactly
\cite[II \S3.4, Proposition~2, II-37--38]{Serre1989}.
\end{proof}

\begin{proposition}\label{prop:frey-weight-open-image}
Let \(\rho:G_\Q\to\operatorname{GL}_2(\Q_\ell)\) be continuous, with
\(\det\rho=\chi_\ell\), and Hodge--Tate at \(\ell\).
Assume that at a density-one set of unramified rational primes \(q\)
its polynomial is \(T^2-a_qT+q\), with \(a_q\in\Z\) and both complex
roots of absolute value \(q^{1/2}\).
If, for one \(r\ne\ell\),
\[
 (\rho|_{G_{\Q_r}})^{\mathrm{ss}}\simeq\delta\chi_\ell\oplus\delta
 \quad\text{with }\delta^2=1,
\]
then \(\rho\) has open image.
\end{proposition}
\begin{proof}
Let \(G\) be the algebraic monodromy group of \(\rho^{\mathrm{ss}}\).
Its connected component is reductive: the fixed space of its unipotent
radical is nonzero and stable under \(G\), and semisimplicity makes
the radical act trivially on every simple summand.
If \(G^\circ\) were a torus, after a finite Galois extension \(L/\Q\)
and coefficient extension its representation would be
\(\psi_1\oplus\psi_2\).
Enlarge \(L\) by a quadratic extension of \(\Q\) killing \(\delta\)
at \(r\). Such an extension is obtained by choosing a rational
representative of the local squareclass.
At each \(w\mid r\), the two local characters are now
\(\chi_\ell\) and \(1\), so both are unramified there.
Choose their moduli without these places.

Lemma~\ref{lem:frey-rank-one-weights} gives constant weights \(b_1,b_2\).
Choose a good rational prime \(q\) splitting completely in \(L\) and
avoiding the finite conductors. Such a prime exists by Chebotarev and
the density-one hypothesis.
Both character values at \(v\mid q\) are roots of
\(T^2-a_qT+q\). Under every complex embedding they remain roots of
this rational polynomial. Thus \(b_1=b_2=1\).
At \(w\mid r\), of residue degree \(f\), the local character values
are \(1\) and \(r^f\). Their absolute values contradict the requirement
that both be \(r^{f/2}\).

Therefore \(G^\circ\) is not a torus. A connected reductive nontorus
subgroup of \(\operatorname{GL}_2\) contains \(\operatorname{SL}_2\); the cyclotomic determinant
then gives \(G^\circ=\operatorname{GL}_2\).
The semisimplification is absolutely irreducible, so \(\rho\) itself
is absolutely irreducible.

For \(H=\rho(G_\Q)\), its Lie algebra \(\mathfrak h\subseteq
\mathfrak{gl}_2\) is \(\operatorname{Ad}(H)\)-stable and therefore
\(\operatorname{Ad}(\operatorname{GL}_2)\)-stable. It is an ideal with nonzero trace because the
determinant image is open. Such an ideal is the center or all of
\(\mathfrak{gl}_2\).
The center would make an open subgroup of \(H\) scalar, contradicting
its full algebraic monodromy. Thus \(\mathfrak h=\mathfrak{gl}_2\),
and \(H\) is open. This Lie argument also works at \(2\), by taking
logarithms on a sufficiently deep congruence subgroup.
\end{proof}

\paragraph{Local geometry and transfer between coefficient primes.}

\begin{lemma}\label{lem:frey-local-two}
For odd \(n\ge5\), every local point or cusp witness at \(2\) gives
\[
 (\rho_{s,3}|_{G_{\Q_2}})^{\mathrm{ss}}
 \simeq\delta\chi_3\oplus\delta,\qquad \delta^2=1.
\]
At either coefficient prime \(\ell=2,3\), the local Frey representation
is potentially semistable.
\end{lemma}
\begin{proof}
For an interior point at \(2\), normalize its projective coordinates
to be integral with at least one unit. Exactly two are units.
After an anharmonic substitution preserving \(j\), the parameter
has valuation \(nm\), with \(m\ge1\). Consequently
\[
 v_2(j)=v_2\!\left(
 2^8\frac{(\lambda^2-\lambda+1)^3}{\lambda^2(\lambda-1)^2}
 \right)=8-2nm<0.
\]
This calculation for prime exponent is already
\cite[Lemmas~5.11--5.12, pp.~21--22]{Steklov2026}; the same calculation
applies to every odd \(n\ge5\).
Negative \(j\)-valuation gives a Tate curve after at most a quadratic
extension \cite[IV Appendix A.1.1--A.1.2, IV-29--32]{Serre1989}.
Its local semisimplification has the asserted form. At a cusp the
Tate description gives the same conclusion.

At a coefficient prime an interior point gives the Tate module of a
local elliptic curve, which is potentially semistable.
A cusp representation is a finite quadratic twist of an extension of
\(\Q_\ell\) by \(\Q_\ell(1)\).
Every such extension is semistable, including at \(\ell=2\), by
\cite[Lemma~8.3.9, p.~125]{BrinonConrad2009}.
One can also see this by spanning \(H^1(K,\Q_\ell(1))\) with Kummer
classes of a uniformizer and of its products with a basis of principal
units. They are Tate-curve classes; sums and scalar multiples are
obtained through direct sums, subobjects and quotients.
\end{proof}

\begin{lemma}\label{lem:frey-companion-comparison}
Suppose continuous representations \(\rho_2\) and \(\rho_3\) of
\(G_\Q\), over \(\Q_2\) and \(\Q_3\), have equal rational Frobenius
polynomials at a density-one set \(\mathcal P\) of primes where both
are unramified. Assume \(\rho_2\) is Hodge--Tate at \(2\), and
\(\rho_3\) has open image. Then the connected algebraic monodromy
of \(\rho_2^{\mathrm{ss}}\) is not a torus.
\end{lemma}
\begin{proof}
If it were a torus, choose a finite Galois \(L/\Q\) and a finite
coefficient extension splitting
\(\rho_2^{\mathrm{ss}}|_{G_L}=\psi_{1,2}\oplus\psi_{2,2}\).
Lemma~\ref{lem:frey-rank-one-weights} supplies algebraic characters over a number
field \(F\) and their realizations \(\psi_{1,3},\psi_{2,3}\), for any
fixed embedding \(F\hookrightarrow\overline\Q_3\).

The prime ideals of \(L\) lying above \(q\in\mathcal P\) which split
completely in \(L\), after finitely many exclusions, have density one
among prime ideals of \(L\).
Indeed, prime ideals of residue degree at least two have a convergent
prime sum at \(s=1\); all unramified degree-one primes lie above
completely split rational primes; and the prime sum above
\(q\notin\mathcal P\) is at most \([L:\Q]\sum_{q\notin\mathcal P}q^{-s}\).
At this density-one set, the Frobenius polynomials of
\(\rho_3|_{G_L}\) and \(\psi_{1,3}\oplus\psi_{2,3}\) agree, because the
algebraic polynomial over \(F\) maps to the given rational polynomial
under the injective two-adic embedding. It therefore equals that
polynomial already in \(F[T]\), and is unchanged by the three-adic embedding.

Their traces agree on all of \(G_L\).
Otherwise, a finite congruence quotient detects a trace difference on
a conjugacy-invariant open subset, to which finite-level Chebotarev
assigns positive density. It must meet the preceding density-one set,
a contradiction.
This argument only discards the finite ramification set of the chosen
finite quotient; it assumes no finite ramification set for \(\rho_3\).
Brauer--Nesbitt gives
\[
 (\rho_3|_{G_L})^{\mathrm{ss}}\simeq\psi_{1,3}\oplus\psi_{2,3}.
\]
Restriction to \(G_L\) preserves openness of the compact image.
An open subgroup has full algebraic monodromy and acts absolutely
irreducibly, contradicting this decomposition.
\end{proof}

\begin{proof}[Proof of Theorem~\ref{thm:prosolvable-frey-open-intro}]
The first assertion is Proposition~\ref{prop:frey-three-mod-nine}.
For a noncuspidal section Lemma~\ref{lem:frey-common-frobenius} gives the common
density-one pure rational polynomials.
Lemmas~\ref{lem:frey-local-two} and Proposition~\ref{prop:frey-weight-open-image},
with \(\ell=3,r=2\), make \(\rho_{s,3}\) open.
Lemma~\ref{lem:frey-companion-comparison} then excludes toroidal connected
monodromy for \(\rho_{s,2}^{\mathrm{ss}}\).
The determinant and Lie-algebra argument in
Proposition~\ref{prop:frey-weight-open-image} makes \(\rho_{s,2}\) open as well.

Let \(H_\ell\) be the individual images and \(H\subseteq H_2\times H_3\)
the joint image. Compact Goursat gives a common profinite quotient.
The images in that quotient of open pro-2 and open pro-3 subgroups
are open. Their intersection is both pro-2 and pro-3, hence trivial.
The quotient is finite, so \(H\) contains a product of open subgroups.
After any finite field extension the individual images remain open,
hence absolutely irreducible with algebraic monodromy \(\operatorname{GL}_2\).
\end{proof}

\paragraph{The remaining arithmetic question.}

The theorem removes potentially abelian semisimplifications, including
reducible, finite-projective, toroidal and dihedral cases, from the
possible noncuspidal quotient sections. It does not prove finite
ramification or modularity in the remaining simultaneous open-image
case.

Khare--Larsen--Ramakrishna construct density-one compatible pure rational
pairs, potentially semistable at their coefficient primes, with infinite
ramification permitted
\cite[\S2.4, Corollary~14, manuscript pp.~23--24]{KhareLarsenRamakrishna2005}.
Their hypotheses require \(p>q\ge5\). This shows why two coefficients
alone do not give a general finite-ramification theorem; it is neither
a counterexample at coefficients \(2,3\) nor a section counterexample.
The usual finite exceptional set in a compatible system cannot be
introduced merely by renaming density-one compatibility.

